\section{Threat Model}\label{sec:threat-model}
With the context from the preceding section, we now begin our study of the security of message deletion in E2EE applications. We start by introducing our threat model in this section. We describe the capabilities and assumptions we make of adversaries, and motivate their goals based on real-world threats and adversarial scenarios. 

\subsection{Overview and Motivation}
The general problem of secure deletion has a long history~\cite{naor2001anti,reardon2013sok,garg2020formalizing,cohn2020delf}, across contexts different from ours. Broadly, the goal of secure deletion is to ``leave no trace''~\cite{garg2020formalizing}: deletion should leave a system in the same state as if the deleted content was not present to begin with, even if an adversary has access to a ``physical medium that contained some representation of the data objects''~\cite{reardon2013sok}.

Carrying over these ideas to our setting, the security goal for E2EE message deletion is that, once deleted, messages should not be recoverable even if given \emph{full, physical access} to a user's device. That is, deleted messages should not be retrieval even by adversaries that are \emph{not} UI-bound~\cite{freed2018stalker}, and can thus inspect the internals of the device: persistent files in storage, local databases, client data structures, etc. In other words, whatever one can learn from the application's \emph{logical} view (i.e., what is displayed in the UI) should be equivalent to what one can from its \emph{physical} view (i.e., the device's storage).

Referring to our systematization from Section~\ref{sec:e2ee-deletion-systematization}, our threat model reflects the guarantee that, after a message deletion feature is triggered by any endpoint (property \#1), an adversary that has access to any device within the feature's deletion scope (property \#2) any time after the message is removed (property \#4) should not be able to learn any information beyond what is revealed by the feature's visible information (property \#3). For instance, if a message is unsent, an adversary with access to the sender or recipient's device (and, thus, their local database of messages) should only be able to learn that a specific message was removed, but not its content.

\paragraph{Motivating scenarios.} There are many natural scenarios in which one would want such privacy guarantees for message deletion features, and for them not to just as UI or storage cleanup. 

For example, a sender may wish to remove an \emph{accidental message} before it is seen by the recipient: perhaps they sent it to the wrong person, or it contains a mistake, etc. As a second example, even if the messages are sent as intended, users may rely on deletion features to make \emph{sensitive messages} ephemeral and safeguard against future device accesses. This could involve exchange of private data (e.g., passwords or credit card numbers), information that could be legally compromising (e.g., whistleblowers, protest organizers, etc.), or discussing confidential information; for instance, in 2025, members of the U.S. government used Signal's disappearing messages feature for military planning~\cite{enwiki:1346011653}.

Perhaps most importantly, users may rely on deletion features to protect against \emph{device coercion}. For example, this may be in settings where the adversary is close to the victim (i.e., \emph{known-adversary} threat models~\cite{}), such in intimate partner violence (IPV)~\cite{} or between family members~\cite{}. Or to protect against compelled disclosure from law enforcement. (threat model section here: \url{https://www.usenix.org/system/files/usenixsecurity23-bellini.pdf})
 
A priori, it may be not obvious why secure deletion is even a relevant goal for E2EE messaging to begin with: If someone has access to a device, can they not just observe all messages to begin with, rendering deletion security moot?  As the examples above show, however, our security goals are reflective of real-world concerns and threats. In particular, it is motivated by the fact that an adversary may have 

Indeed, this matches perceptions and expectations of users. \andres{Cite literature on user perceptions.}

Indeed [say apps offer this feature for a reason and cite. And that it is under privacy menus.]

\subsection{Threat Model Details}
Let us now unpack our threat model in detail. To do so, we first overview the general architecture of E2EE messaging applications, which will allow us to explain how our threat model slots in this broader system.
\paragraph{Architecture of E2EE applications.} We depict a typical architecture for an E2EE messaging application in Figure~\ref{}. (Note that we purposefully omit many details, focusing on what is most relevant to our threat model.) In this setting, users exchange messages via service-operated servers, with the key security goal being that not even a malicious or compromised server can read nor modify the plaintext messages. To achieve this, messages are routed using an E2EE protocol (e.g., the Signal protocol~\cite{signal-docs}), which encrypts and decrypts messages at the endpoints with keys that are not known by the server. Users participate in this protocol via \emph{application clients} installed in their devices, which abstract internal application operations (e.g., cryptographic operations) behind the client's user interface (UI). Notice that E2EE security thus (implicitly or explicitly) assumes that client software is trusted~\cite{fabrega2025mitigating}.

The client stores all relevant user information (such as sent or received messages, local settings, etc.) in the device's internal storage; we refer to this stored information as the user's \emph{application state}, and denote it by $\appState$. For example, for Android applications, a common structure for $\appState$ consists of a primary SQLite database (serialized in some way) for most application information, and an auxiliary BLOB store for larger binary files (e.g., exchanged media)~\cite{spahn2014pebbles,fabrega2025mitigating}. The client communicates with the device's storage through some model-dependent API (e.g., SQL commands for SQLite) to read from and write to $\appState$. Individual application-level operations often result in various storage-level transactions from the client to $\appState$. For example, sending a message may simultaneously create a new row in a table of messages, modify a row in a table of active conversations, update numerous local variables, etc.

To support various features, the client may periodically \emph{export} (a function of) $\appState$ outside the user's device. For example, the application may support E2EE backups~\cite{signal-e2ee-backups,whatsapp-e2ee-backups}, which store an encrypted copy of $\appState$ in some cloud, allowing users to restore their information (e.g., chat history) should they lose access to their device.\andres{Probably remove.}

\paragraph{Overview of threat model.}Let us now explain our threat model. For clarity of presentation, we keep our discussion here conceptual, in order to focus on the main ideas. Looking ahead, we will provide a formal definition of security in Section~\ref{sec:security-definition}.

\andres{Maybe move intuition before architecture? And then say let us unpack this.}
If an E2EE application implements message deletion securely, then deleted messages (which are not present in the UI) should not be recoverable from the device's physical state. If, however, the implementation is not secure, straggling fragments of deleted messages may be recoverable from having direct access to storage.\andres{Clarify that, as we will see later on, this may arise from straggling deletion (a la forensics), or from subtle application logic such as deduplication.}


Let us unpack this in more detail. At a high level, our threat model is underpinned by one main ingredient: We consider an adversary that has \emph{gained access to a user's local application state $\appState$}, from which it tries to recover information about messages that were deleted prior to their access. 
% In particular, this means the adversary has access to the client's \emph{physical} state, such as the serialized database file in storage.
That is, let $t_0$ be the time when a message deletion feature is triggered and $t_1 > t_0$ be the time when some message $M$ is deleted (as mentioned in Section~\ref{sec:e2ee-deletion-systematization}, different features may incur different time delays). The adversary gains direct access to $\appState$ at some time $t_2 > t_1$. By inspecting $\appState$, the adversary tries to learn some information about $M$.

Access to $\appState$ can arise from, most notably, access to the user's device; or from access to some external location where $\appState$ has been exported to.\footnote{The latter may need to be coupled with device access as well, in order to recover the secret key used to encrypt the application state.} In practice, the way in which one can obtain $\appState$ depends on the specific platform and application in question; we show a few examples in Section~\ref{sec:attacks}.

\paragraph{Out-of-scope attacks.} We explicitly leave some attacks outside our threat model.

First, we assume the adversary does not tamper with client code. As mentioned earlier, this is inherent to all E2EE threat models; if client code is not trusted, all confidentiality guarantees are lost to begin with (e.g., an adversary can install a backdoor that directly leaks all sent and received messages).

Second, to keep our scope focused, we limit attention to \emph{message} deletion. One could also consider \emph{metadata} deletion, for instance, that it should not be possible to know that two parties communicated after their chat has been closed. We leave this as an interesting direction for future work.

Lastly, we stress that, in this work, we consider adversaries at the \emph{application level}, and not the \emph{cryptographic-protocol level}. Deletion at the protocol level is an orthogonal threat model, which bears relationship with well-studied security notions such as deniable encryption~\cite{} and post-compromise security~\cite{}.


% \paragraph{E2EE applications} We depict a typical architecture for an E2EE application in Figure~\ref{}, which we overview below.

% We denote by $\user$ a user of an E2EE application service $\service$. The user owns a device with an application client for $\service$ installed, denoted by $\client$, which allow $\user$ to exchange content with other users of the application. This content is relayed as encrypted messages through service-operated servers, in accordance with some application-specific E2EE protocol, such as the Signal protocol~\cite{signal-doubleratchet,signal-x3dh}. 

% The client $\client$ stores all relevant user information (e.g., sent and received messages, local settings, etc.) in the device's internal storage; we refer to this stored information as $\user$'s \emph{application state}, and denote it by $\state$. For example, for Android applications, a common structure for $\state$ consists of a primary SQLite database for most application information, and an auxuliary BLOB store for larger binary files (e.g., exchanged images)~\cite{spahn2014pebbles,fabrega2025mitigating}. The structure of $\state$ is governed by the application's \emph{data model}, or schema, which defines the information that is stored and how it is arranged. For example, for an SQLite database, the data model defines all tables and columns of the database.


% The client $\client$ communicates with the device's storage through some model-dependent API (e.g., SQL commands) to read from and write to $\state$. Individual application-level operations often result in various storage-level transactions from $\client$ to $\state$. For example, sending a message may simultaneously create a new row in a $\messagesTable$, update rows in $\recipientsTable$ and $\conversationsTable$ tables, etc.

% To support various features, $\client$ may periodically \emph{export} $\state$ outside $\user$'s device. For example, the application may support backups, which requires storing a snapshot of $\state$ in some cloud. Or the application may allow for support multiple user devices, which requires propagating $\state$ from $\client$ to a new device whenever $\user$ registers one. Such exports generally involve $\client$ encrypting $\state$ locally using an application-specific symmetric encryption scheme and a per-user secret key (which is sometimes derived from a user password), and uploading the resulting ciphertext to some location outside the device (the cloud, another of $\user$'s devices, etc).

